% TOP_PROBLEM: {"id": 30007028, "title": "Generalized lc flip termination on Kähler fourfolds", "category_id": 5, "category": {"id": 5, "name": "algebraic_geometry", "display_name": "Algebraic Geometry", "description": "Geometric objects defined by polynomial equations.", "slug": "algebraic-geometry", "order_index": 5, "created_at": "2026-07-31T15:26:25.671Z"}, "status": "open"}
% FIRST_POSTED: 2026-10-06
% ATTEMPT_STATUS: solved
% OPENAI_CLAIM: {"schema_version":1,"source_repo":"https://github.com/openai/math","source_commit":"adc7f1241b42e322a6451854ab7e4b4c146bf78a","release_date":"2026-10-06","match":"full","resolution":"proved","independently_reviewed":false,"second_pass":"agreed","families":[{"family":"056","title":"Termination of projective and Kähler fourfold minimal model programs","subject":"Algebraic and complex geometry","manuscripts":[{"title":"Termination of generalized log canonical flips on compact Kähler fourfolds","pdf_path":"preprints/Termination-of-generalized-log-canonical-flips-on-compact-Kahler-fourfolds-October-5-2026/termination-generalized-lc-kahler-fourfolds.pdf","readme_path":"preprints/Termination-of-generalized-log-canonical-flips-on-compact-Kahler-fourfolds-October-5-2026/README.md","date":"2026-10-05","theorem_ref":"thm:main (sections/01-introduction.tex:25)","statement_sources":[{"path":"preprints/Termination-of-generalized-log-canonical-flips-on-compact-Kahler-fourfolds-October-5-2026/build/sections/01-introduction.tex","line":25,"label":"thm:main"}]}],"lean":null,"reasoning_trace":null}]}
\documentclass[11pt]{article}
\usepackage[margin=25mm]{geometry}
\usepackage{fontspec}
\setmainfont{Latin Modern Roman}
\usepackage{amsmath,amssymb,mathtools}
\usepackage{unicode-math}
\setmathfont{Latin Modern Math}
\usepackage{enumitem}
\usepackage{xurl,hyperref,bookmark}
\hypersetup{colorlinks=true,urlcolor=blue}
\setcounter{secnumdepth}{0}
\setlength{\parindent}{0pt}
\setlength{\parskip}{5pt}
\setlength{\emergencystretch}{3em}
\begin{document}
\section*{Generalized lc flip termination on Kähler fourfolds}

\subsection{OpenAI's claim}
Every existing generalized log-canonical flip sequence in the specified globally factorial compact Kähler fourfold category is claimed to terminate, without scaling or pseudo-effectivity assumptions. The cited manuscript argues for that result; its relation to this problem is specified below.

\subsection{Mathematical statement of the claimed result}
\paragraph{Mathematical claim 1.} Let $X_0$ be normal irreducible globally Weil rational-factorial compact Kähler of dimension four, with effective rational finite-support $B_0$ of coefficients in $[0,1]$ and a fixed analytically nef rational b-divisor represented on a projective bimeromorphic compact Kähler model. Suppose $D_0=K_{X_0}+B_0+M_0$ rational Cartier and all generalized log discrepancies nonnegative. With strict transforms, the same b-divisor and compatible canonical representatives, OpenAI claims finiteness of every sequence
\[
X_i\xrightarrow{f_i}Z_i\xleftarrow{f_i^+}X_{i+1},\qquad
-D_i\text{ }f_i\text{-ample},\quad D_{i+1}\text{ }f_i^+\text{-ample},
\]
with every $Z_i$ normal compact Kähler, whose maps are projective small nonisomorphic bimeromorphic maps with connected fibers and whose models remain normal irreducible globally Weil rational-factorial compact Kähler fourfolds with rational Cartier adjoints. The coefficients and discrepancy bound permit generalized log canonical, not only generalized canonical, pairs.

\subsection{Outline of the argument}
Manuscript 1 (family 056, 2026-10-05): Relative programs and special termination reduce the sequence to controlled terminal pieces while preserving the nef datum and actual adjoints. Analytic cycle ranks remove divisorial losses, and strict mixed-Hodge weight comparisons bound the residual surface changes, leaving only finitely many possible steps.

\subsection{Scope relative to this problem}
Manuscript 1 (family 056, 2026-10-05): The problem selects the mathematical statement summarized here, including its stated hypotheses and parameter range. The cited theorem gives a claimed proof of that assertion in full.

\subsection{Formal verification}
Family 056 has no Lean scope page in this snapshot; no formalization of these manuscript results is listed.

This site has not executed the Lean files or checked their dependencies. These descriptions report the scope of the linked declarations, rather than a new verification of the proofs.

\subsection{Status caveat}
OpenAI attributes these results to an internal model. Its release notice cautions that results lacking full formalization can contain mistakes. This site has not independently assessed the proofs; a solved label records the scope of the reported claim.

\subsection{Sources}
\begin{itemize}
\item OpenAI, \emph{Termination of generalized log canonical flips on compact Kähler fourfolds}, 2026-10-05; citation key \texttt{OAI:Termination-of-generalized-log-canonical-flips-on-compact-Kahler-fourfolds-October-5-2026}. \href{https://github.com/openai/math/blob/adc7f1241b42e322a6451854ab7e4b4c146bf78a/preprints/Termination-of-generalized-log-canonical-flips-on-compact-Kahler-fourfolds-October-5-2026/termination-generalized-lc-kahler-fourfolds.pdf}{PDF (GitHub)}; \href{https://github.com/openai/math/raw/adc7f1241b42e322a6451854ab7e4b4c146bf78a/preprints/Termination-of-generalized-log-canonical-flips-on-compact-Kahler-fourfolds-October-5-2026/termination-generalized-lc-kahler-fourfolds.pdf}{PDF (direct)}; \href{https://github.com/openai/math/blob/adc7f1241b42e322a6451854ab7e4b4c146bf78a/preprints/Termination-of-generalized-log-canonical-flips-on-compact-Kahler-fourfolds-October-5-2026/README.md}{Manuscript page and citation}; \href{https://github.com/openai/math/blob/main/preprints/Termination-of-generalized-log-canonical-flips-on-compact-Kahler-fourfolds-October-5-2026/termination-generalized-lc-kahler-fourfolds.pdf}{Latest version}; Statement source: \href{https://github.com/openai/math/blob/adc7f1241b42e322a6451854ab7e4b4c146bf78a/preprints/Termination-of-generalized-log-canonical-flips-on-compact-Kahler-fourfolds-October-5-2026/build/sections/01-introduction.tex\#L25}{thm:main}.
\item \href{https://github.com/openai/math/blob/adc7f1241b42e322a6451854ab7e4b4c146bf78a/CONTENTS.md}{OpenAI catalogue entry}.
\end{itemize}
\end{document}
